\documentclass[12pt,a4paper]{report}
\usepackage{fontspec}
\setmainfont{Noto Serif}
\setmonofont{Noto Sans Mono}
\usepackage{polyglossia}
\newfontfamily\cyrillicfonttt{Noto Sans Mono}
\newfontfamily\titlepagefont{Noto Sans}
\setdefaultlanguage{mongolian}
\usepackage[a4paper,margin=2.5cm]{geometry}
\usepackage{setspace,graphicx,float,booktabs,longtable,array,hyperref,xurl,enumitem,listings,amssymb,xcolor,ragged2e}
\newcolumntype{P}[1]{>{\RaggedRight\hspace{0pt}}p{#1}}
\onehalfspacing
\hypersetup{colorlinks=true,linkcolor=black,urlcolor=blue,pdftitle={Олон урсгалт Sudoku шалгах програмын тайлан}}
\setcounter{tocdepth}{2}
\graphicspath{{docs/screenshots/}{./}}

% ---- Макронууд ----
\newcommand{\repo}{\url{https://github.com/Naidan2006/sudoku-validator}}
\newcommand{\shot}[3]{\begin{figure}[H]\centering\includegraphics[width=0.85\textwidth]{#1}\caption{#2}\label{#3}\end{figure}}
% 9x9 хүснэгт: 9 мөр, мөр бүр "a&b&c&d&e&f&g&h&i" хэлбэртэй
\newcommand{\sgrid}[9]{{\small\setlength{\tabcolsep}{3.5pt}\begin{tabular}{|ccc|ccc|ccc|}\hline #1\\#2\\#3\\\hline #4\\#5\\#6\\\hline #7\\#8\\#9\\\hline\end{tabular}}}

\lstset{basicstyle=\ttfamily\small,breaklines=true,frame=single}

\begin{document}
%==================== НҮҮР ====================
\begin{titlepage}
\vspace*{0.5cm}
\begin{center}
{\titlepagefont\large ШИНЖЛЭХ УХААН ТЕХНОЛОГИЙН ИХ СУРГУУЛЬ}\par\vspace{0.3cm}
{\titlepagefont\normalsize МЭДЭЭЛЭЛ, ХОЛБООНЫ ТЕХНОЛОГИЙН СУРГУУЛЬ}\par\vspace{1.5cm}
\IfFileExists{logo.png}{\includegraphics[width=3.5cm]{logo.png}}{\fbox{\parbox[c][3.5cm][c]{3.5cm}{\centering logo.png}}}\par\vspace{1.5cm}
{\titlepagefont\Large F.CSM302 -- ҮЙЛДЛИЙН СИСТЕМ}\par\vspace{0.7cm}
{\titlepagefont\Large БИЕ ДААЛТ 1}\par\vspace{1.2cm}
{\titlepagefont\LARGE\textbf{SUDOKU SOLUTION VALIDATOR}}\par\vspace{0.4cm}
{\titlepagefont\large Олон урсгалт Sudoku шийдэл шалгах програм}\par\vspace{0.4cm}
{\titlepagefont\large Эцсийн тайлан}\par\vspace{1.2cm}
\begin{minipage}{0.92\textwidth}
\titlepagefont\normalsize
\textbf{Хуваарилалтын код:} II-4-5\par\vspace{0.25cm}
\textbf{Сурах бичгийн төсөл:} Operating System Concepts (2018), Chapter 4: Threads \& Concurrency, Programming Project 5\par\vspace{0.25cm}
\textbf{Гүйцэтгэсэн багийн гишүүд:}\par\vspace{0.15cm}
Д.~Алтан-Очир (B221940001)\par
Б.~Аззаяа (B242270803)\par
О.~Энхжаргал (B232230049)\par
Н.~Мөнхбаяр (B242270058)\par
Ц.~Тэргэл (B241940128)\par
А.~Найдан (B241930004)\par\vspace{0.3cm}
\textbf{Шалгасан багш:} Г. Ганбат /F.CS10/\par\vspace{0.25cm}
\textbf{GitHub:} \url{https://github.com/Naidan2006/sudoku-validator}
\end{minipage}
\vfill
{\titlepagefont Улаанбаатар хот\\2026 он}
\end{center}
\end{titlepage}
\pagenumbering{roman}\tableofcontents\clearpage\pagenumbering{arabic}

%==================== ХУРААНГУЙ ====================
\chapter*{Хураангуй}\addcontentsline{toc}{chapter}{Хураангуй}
\textbf{Асуудал.} Sudoku-ийн бөглөсөн 9$\times$9 хүснэгт зөв эсэхийг мөр, багана, 3$\times$3 дэд хүснэгт бүрээр шалгах шаардлагатай бөгөөд эдгээр шалгалт нь бие даасан тул зэрэгцүүлэн гүйцэтгэхэд тохиромжтой.

\textbf{Шийдэл.} C хэл болон POSIX Threads (\texttt{pthread}) ашиглан шалгалтыг нийт 11 урсгалд хуваасан: мөрүүдийг шалгах 1, багануудыг шалгах 1, есөн дэд хүснэгтийг тус бүр шалгах 9 урсгал. Үндсэн урсгал тэдгээрийг \texttt{pthread\_create()}-ээр үүсгэж, \texttt{pthread\_join()}-ээр хүлээж, үр дүнг нэгтгэн VALID/INVALID гэж гаргана. GTK3 график интерфэйс нь хэрэглэгчид хүснэгтийн утгыг оруулах, шалгах, засаж дахин шалгах, цэвэрлэх боломж олгож, 11 урсгал тус бүрийн төлөвийг (WAITING, RUNNING, PASS, FAIL, ERROR) тусгай самбарт харуулна.

\textbf{Туршилт.} Хоосон нүд, буруу Sudoku, мөр/багана/дэд хүснэгтийн тусгаарлагдсан алдаа, буруу тэмдэгт, засварлан дахин шалгах, цэвэрлэх зэрэг 10 туршилтын кейсийг боловсруулж, оролтын хүснэгтийг тайланд бүрэн оруулсан. Гол үр дүн: програм хоосон нүдийг мөр, баганын байрлалтай нь мэдэгдэж, алдаатай хэсгийг FAIL, зөв хүснэгтийг бүх 11 хэсэг PASS болгож VALID гэж үнэлдэг.\\[0.3em]
\textbf{Түлхүүр үг:} үйлдлийн систем, олон урсгал, pthread, Sudoku, GTK3.

%==================== 1 ====================
\chapter{Удиртгал}
Олон урсгалт програмчлал нь нэг програмын дотор бие даасан ажлуудыг зэрэгцүүлэн гүйцэтгэх боломж олгодог. Энэхүү төсөл нь үйлдлийн системийн урсгал үүсгэх, дуусахыг хүлээх болон үр дүнг нэгтгэх ойлголтыг бодит жишээгээр хэрэгжүүлсэн. Sudoku нь 9$\times$9 хүснэгтэд 1-ээс 9 хүртэлх тоог мөр, багана, 3$\times$3 дэд хүснэгт бүрд давхардалгүй байрлуулах дүрэмтэй.
\section{Төслийн зорилго}
C хэл, POSIX Threads ашиглан Sudoku хүснэгтийн мөр, багана, дэд хүснэгтийг олон урсгалаар шалгадаг, үр дүнг ойлгомжтой харуулдаг програм боловсруулах.
\section{Зорилтууд}
\begin{itemize}
\item 9$\times$9 хүснэгтийн өгөгдлийг хэрэглэгчээс хүлээн авах (эх кодод хатуу тогтоохгүй).
\item Мөр, багана, есөн дэд хүснэгтийн шалгалтыг тус тусад нь гүйцэтгэх.
\item \texttt{pthread\_create()} болон \texttt{pthread\_join()} ашиглан урсгалуудыг зохицуулах.
\item Хоосон нүд, буруу оролтод ойлгомжтой мэдээлэл өгөх.
\item GTK3 интерфэйсээр үр дүнг харуулж, засварлан дахин шалгах боломж олгох.
\item Зөв болон буруу өгөгдлөөр туршилт хийж, үр дүнг дахин шалгаж болохуйц баримтжуулах.
\end{itemize}
\section{Ач холбогдол}
Шалгалтын хэсгүүд хоорондоо хамааралгүй тул олон урсгалд хуваахад тохиромжтой бөгөөд урсгал үүсгэх, хүлээх, үр дүн нэгтгэх, UI-г гацаахгүй байх зэрэг үйлдлийн системийн үндсэн ойлголтыг нэг жижиг төсөл дээр практикаар хэрэгжүүлэх боломж олгодог.
\section{Хамрах хүрээ}
Програм нь Sudoku-г бодож шийдэхгүй, хэрэглэгчийн өгсөн 9$\times$9 хүснэгтийн зөв эсэхийг шалгана. Мөр бүр, багана бүр болон есөн 3$\times$3 дэд хүснэгтэд 1--9 тоо тус бүр нэг удаа орсон эсэхийг шалгана. Консол болон GTK3 график гэсэн хоёр хувилбартай.

%==================== 2 ====================
\chapter{Сурах бичгийн шаардлага}
Энэхүү төсөл нь Silberschatz, Galvin, Gagne нарын \textit{Operating System Concepts} (2018) сурах бичгийн 4-р бүлэг \textit{Threads \& Concurrency}-ийн Programming Project 5 буюу Sudoku Solution Validator юм. Хуваарилалтын код: \textbf{II-4-5}. Сурах бичгийн төсөл нь 9$\times$9 Sudoku шийдлийг шалгахдаа бүх мөрийг шалгах нэг урсгал, бүх баганыг шалгах нэг урсгал, есөн 3$\times$3 дэд хүснэгтийг тус бүр шалгах есөн урсгал, нийт 11 урсгал ашиглахыг шаарддаг. Доорх хүснэгтэд шаардлага ба хэрэгжилтийн холбоог харуулав.
\begin{table}[H]\centering\caption{Шаардлага ба хэрэгжилт}\begin{tabular}{P{0.4\textwidth}P{0.5\textwidth}}\toprule
\textbf{Шаардлага}&\textbf{Хэрэгжилт}\\\midrule
Үндсэн функц&Мөр, багана, 3$\times$3 дэд хүснэгтийг шалгана.\\
Олон урсгал&Мөр шалгах 1, багана шалгах 1, дэд хүснэгт шалгах 9 урсгал.\\
Өөр оролтоор турших&Хэрэглэгч утгаа оруулж, засаж, дахин шалгана; цэвэрлэх товчтой.\\
Алдааны мэдээлэл&Хоосон нүд, буруу тэмдэгт, дүрэм зөрчсөн өгөгдлийг илрүүлнэ.\\
UI дүрслэл&GTK3 интерфэйс PASS/FAIL, VALID/INVALID болон 11 урсгалын төлөв харуулна.\\
Ажиллуулах заавар&GitHub README болон Manual.pdf-д заавар байршуулсан.\\\bottomrule\end{tabular}\end{table}

%==================== 3 ====================
\chapter{Онолын үндэс}
\section{Процесс ба урсгал}
Процесс нь програм ажиллах үеийн нөөц, санах ойн орчныг агуулна. Урсгал (thread) нь процесс дотор гүйцэтгэгдэх ажлын нэг дараалал юм. Нэг процессын урсгалууд өгөгдөл, кодоо хуваалцдаг тул хамтын өгөгдөлд хандах ажиллагааг зөв зохицуулах шаардлагатай.
\section{POSIX Threads}
\texttt{pthread} нь POSIX орчинд урсгал үүсгэх, удирдах интерфэйс юм. \texttt{pthread\_create()} нь шинэ урсгал үүсгэж, \texttt{pthread\_join()} нь тухайн урсгал дуусахыг хүлээнэ. Програм эдгээр функцийн буцаах утгыг шалгаж, бүх шалгалт дууссаны дараа үр дүнг нэгтгэх ёстой.
\section{Sudoku шалгах алгоритм}
Зөв Sudoku-д мөр бүр, багана бүр болон 3$\times$3 дэд хүснэгт бүр 1--9 тоог тус бүр яг нэг удаа агуулна. Хэсэг бүрийг шалгахдаа 1--9 тоо бүрийн тохиолдох тоог тоолж, аль нэг тоо хоёр ба түүнээс олон удаа эсвэл огт тохиолдохгүй бол тухайн хэсгийг FAIL гэнэ. Бүх хэсэг PASS бол хүснэгт VALID, нэг ч хэсэг FAIL бол INVALID гэж үзнэ.
\section{Урсгалын хуваарилалт ба өгөгдөл хуваалцах}
Програм 11 шалгах урсгал ашиглана: нэг нь мөрүүд, нэг нь баганууд, ес нь тус тусын 3$\times$3 дэд хүснэгтийг шалгана. Урсгал бүр хүснэгтийг зөвхөн унших бөгөөд үр дүнгээ \texttt{results[]} массивын өөрийн тусдаа элементэд бичдэг (мөр: \texttt{results[0]}, багана: \texttt{results[1]}, дэд хүснэгт: \texttt{results[index+2]}) тул хоёр урсгал нэг нүдэнд бичихгүй. Иймээс уралдаан (race condition) үүсэхгүй бөгөөд mutex шаардлагагүй. Үндсэн урсгал бүх урсгалыг \texttt{pthread\_join()}-оор хүлээсний дараа л үр дүнг нэгтгэнэ. Шалгалтын явцад хүснэгтийн өгөгдлийг өөрчлөхгүй байхаар GUI талыг зохион байгуулсан (5.5 хэсэг).

%==================== 4 ====================
\chapter{Системийн зохиомж}
\section{Ерөнхий бүтэц}
\begin{figure}[H]\centering\fbox{\parbox{0.82\textwidth}{\centering Хэрэглэгчийн 9$\times$9 оролт\\$\downarrow$\\GTK3 интерфэйс (\texttt{gui.c})\\$\downarrow$\\Шалгах логик (\texttt{sudoku\_core.c})\\$\downarrow$\\11 pthread урсгал: мөр (1), багана (1), дэд хүснэгт (9)\\$\downarrow$\\Үр дүн нэгтгэж PASS/FAIL, VALID/INVALID харуулах}}\caption{Системийн өгөгдлийн урсгал}\end{figure}
\section{Репозиторийн бүтэц}
\begin{table}[H]\centering\caption{Репозиторийн бүтэц}\begin{tabular}{P{0.38\textwidth}P{0.52\textwidth}}\toprule
\textbf{Зам}&\textbf{Үүрэг}\\\midrule
\texttt{README.md}&Төслийн тайлбар, хурдан ажиллуулах заавар.\\
\texttt{main.c}&Консолын хувилбарын оролт, үр дүнгийн харуулалт.\\
\texttt{gui.c}&GTK3 график интерфэйс, хүснэгтийн утга ба үр дүн.\\
\texttt{sudoku\_core.c}&Мөр, багана, дэд хүснэгтийн шалгалт, урсгалын зохицуулалт.\\
\texttt{sudoku\_core.h}&Хэмжээ, урсгалын тоо, функцийн зарлалууд.\\
\texttt{tests/}&Тестийн тайлбар (README).\\
\texttt{.github/workflows/build.yml}&GitHub Actions-ийн автомат build.\\
\texttt{docs/}&Гарын авлага (Manual.pdf), дэлгэцийн зургууд.\\
\texttt{tailan\_latex/main.tex}, \texttt{OS\_BD\_1.pdf}&Тайлангийн LaTeX эх ба PDF.\\\bottomrule\end{tabular}\end{table}
Эх код үндсэн хавтаст байрладаг бөгөөд GitHub Actions-ийн build энэ замаар ажилладаг. Тестийн оролтыг тайлангийн 7.2 хэсэгт бүрэн өгч, үр дүнг screenshot-оор баримтжуулсан.
\section{Өгөгдөл ба үр дүнгийн урсгал}
Хэрэглэгч хүснэгтийн утгыг интерфэйсээр оруулна. Програм оролтыг шалгаж (хоосон нүд, нэг оронтой цифр эсэх, 1--9-ийн хооронд эсэх), зөв бол утгыг \texttt{sudoku[row][col]}-д хадгалан шалгах логикт дамжуулна. Урсгалууд өөр өөр хэсгийг шалгаж үр дүнгээ хадгална. Үндсэн урсгал бүх шалгалт дууссаны дараа үр дүнг нэгтгэн харуулна.

%==================== 5 ====================
\chapter{Хэрэгжилт}
\section{Ашигласан технологи}
\begin{itemize}\item C хэл -- шалгалтын алгоритм болон үндсэн логик.\item POSIX Threads (\texttt{pthread}) -- урсгал үүсгэх, дуусахыг хүлээх.\item GTK3 -- график интерфэйс.\item GCC -- эх код хөрвүүлэх.\item Windows болон MSYS2 UCRT64 -- хөгжүүлж, туршсан орчин.\end{itemize}
\section{Орчин ба хувилбар}
\begin{table}[H]\centering\caption{Хөгжүүлэлт ба туршилтын орчин}\begin{tabular}{P{0.35\textwidth}P{0.55\textwidth}}\toprule
\textbf{Бүрэлдэхүүн}&\textbf{Хувилбар}\\\midrule
Үйлдлийн систем&Windows 11 Home, Version 25H2, OS build 26200.9550\\
Орчин&MSYS2 UCRT64\\
GCC&16.2.0 (MSYS2 UCRT64)\\
GTK3&3.24.52\\
POSIX Threads&MSYS2 (winpthreads)\\\bottomrule\end{tabular}\end{table}
\section{Гол функцүүд}
\begin{itemize}\item \texttt{check\_rows()} -- бүх мөрийг шалгана.\item \texttt{check\_columns()} -- бүх баганыг шалгана.\item \texttt{check\_subgrid()} -- нэг 3$\times$3 дэд хүснэгтийг шалгана.\item \texttt{validate\_sudoku()} -- шалгах урсгалуудыг үүсгэж, хүлээж, үр дүнг нэгтгэнэ.\end{itemize}
Урсгал үүсгэх, хүлээх хэсгийн хялбаршуулсан загвар:
\begin{lstlisting}[language=C]
pthread_t tid[TOTAL_THREADS];            /* TOTAL_THREADS = 11 */
/* results[0]: rows, results[1]: columns, results[idx+2]: subgrid idx */
for (int i = 0; i < TOTAL_THREADS; i++) {
    if (pthread_create(&tid[i], NULL, worker, &params[i]) != 0) {
        /* creation failed: join the threads created so far */
    }
}
for (int i = 0; i < TOTAL_THREADS; i++)
    pthread_join(tid[i], NULL);          /* wait for all checks */
/* merge: VALID only if every results[i] is PASS */
\end{lstlisting}
Дээрх нь эх кодын бүтцийг харуулсан хялбаршуулсан загвар бөгөөд \texttt{worker}, \texttt{params} нь жишээний нэр юм.
\section{Хөрвүүлэх ба ажиллуулах}
GTK3 суулгасан MSYS2 UCRT64 терминалд:
\begin{lstlisting}
git clone https://github.com/Naidan2006/sudoku-validator
cd sudoku-validator
gcc gui.c sudoku_core.c -o sudoku_gui.exe -pthread $(pkg-config --cflags --libs gtk+-3.0)
./sudoku_gui.exe
\end{lstlisting}
Консолын хувилбар:
\begin{lstlisting}
gcc main.c sudoku_core.c -o sudoku_validator.exe -pthread
./sudoku_validator.exe
\end{lstlisting}
Консолын хувилбар нь ажиллах үедээ \texttt{Row 1:}--\texttt{Row 9:} гэж асуух бөгөөд хэрэглэгч 9 мөрийн утгыг гараар оруулна.
\section{Хэрэглэгчийн интерфэйс}
График интерфэйс нь Sudoku хүснэгтийн 81 нүд, шалгах, цэвэрлэх (reset), заавар үзэх товч болон шалгалтын үр дүнг харуулах хэсэгтэй. Хэрэглэгч утгаа оруулж шалгаад, үр дүнгээс хамааран өгөгдлөө засаж дахин шалгана.

Интерфэйсийн responsive байдлыг сайжруулахын тулд үндсэн хуудасны агуулгыг босоо гүйлгэх боломжтой болгож, зааврын хуудасны агуулгыг гүйлгэх хэсэгт байрлуулсан. Ингэснээр цонхны өндөр бага үед ч зааврын агуулга гүйлгэгдэж, эхлэх товч доод хэсэгт харагдана. Урсгалын төлөвийн самбарын 11 төлвийг хоёр баганад байрлуулж, дэлгэцийн өргөнд тохируулан текст мөр шилжихээр зохион байгуулсан.

Шалгах товч дарахад \texttt{on\_check\_clicked()} нүд бүрийн утгыг уншиж, хоосон нүд илэрвэл мөр, баганын байрлалтай алдааны мэдэгдэл өгнө; утга нь нэг оронтой цифр биш эсвэл 1--9-ийн гадуур бол алдаа мэдэгдэнэ. Зөв бол шалгалтыг \texttt{g\_thread\_new()}-ээр үүсгэсэн тусдаа GLib worker урсгалд ажиллуулах бөгөөд тэр доторх \texttt{validate\_sudoku()} нь 11 POSIX урсгалыг үүсгэнэ. Ингэснээр шалгалт явагдах үед интерфэйс гацахгүй. Дууссаны дараа үр дүнг \texttt{g\_idle\_add()}-ээр GTK-ийн үндсэн урсгалд шинэчилнэ. Шалгалт явагдаж байх үед цэвэрлэх (Clear) товч блоклогдоно. Урсгалын төлөвийн самбар нь мөр шалгах нэг, багана шалгах нэг, дэд хүснэгт тус бүрийг шалгах есөн урсгалын төлвийг харуулах бөгөөд төлвийн утгууд нь WAITING, RUNNING, PASS, FAIL болон ERROR байна.
\section{GitHub дээрх хамтын ажиллагаа}
Эх код, баримт бичгийг GitHub репозиторид (\repo) байршуулсан. \texttt{main} branch дээр нийт 24 commit байна; өөр нэмэлт branch харагдаагүй. \texttt{git shortlog} нь 20 commit-ийг \texttt{Naidan2006}, 4 commit-ийг \texttt{Naidan} нэрээр харуулсан бөгөөд хоёр нэр ижил email ашигласан тул эдгээр тоо нь багийн зургаан гишүүний commit-ийн тус тусын тоо биш. Хувилбар \texttt{v1.0.0} tag-ийг 2026-10-17-нд \texttt{6b43f4c} commit дээр тавьсан бөгөөд хамгийн сүүлийн commit нь \texttt{4222485}.
\begin{table}[H]\centering\caption{Гол commit-үүд}\begin{tabular}{P{0.2\textwidth}P{0.65\textwidth}}\toprule
\textbf{Commit}&\textbf{Агуулга}\\\midrule
\texttt{6b43f4c}&Тестийн README зөв байршилд нэмсэн (tag \texttt{v1.0.0})\\
\texttt{91a6fc0}&Автомат build workflow\\
\texttt{70851b3}&GUI ба screenshot-ууд\\
\texttt{d4d26a7}&README баримтжуулалт\\
\texttt{43afe01}&GUI, баримтжуулалт, workflow шинэчилсэн\\
\texttt{62594c7}&README-ийн өөрчлөлтүүдийг merge хийсэн\\
\texttt{4222485}&Эцсийн тайлангийн PDF ба LaTeX эх\\\bottomrule\end{tabular}\end{table}
Commit \texttt{62594c7} дээр GitHub Actions-ийн build амжилттай (Success) дууссан. Өөрчлөлтийг \texttt{git push origin main}-аар шууд оруулж, локал дээр \texttt{git merge} хийсэн. Хувилбарыг \texttt{v1.0.0} tag-аар тэмдэглэсэн. Гишүүн бүрийн хариуцсан ажлыг Хавсралт А-д үзүүлэв.
\section{Бэлэн код, AI-ийн хэрэглээ}
Төслийн алгоритм ба урсгалын зохион байгуулалтыг сурах бичгийн шаардлагын дагуу багийнхан хэрэгжүүлсэн бөгөөд гадны сан болгон зөвхөн GTK3 болон pthread-ийг ашигласан. Дараах зорилгоор AI хэрэгсэл ашигласан:
\begin{itemize}
\item ChatGPT -- код, тайлбар, баримтжуулалт боловсруулахад туслах.
\item Claude -- LaTeX тайлангийн алдаа засах, бүтэц, агуулгыг зааврын шаардлагатай тулган сайжруулахад.
\end{itemize}
AI-ийн гаргасан код, текстийг багийн гишүүд хянаж, хөрвүүлж ажиллуулан туршсан.

%==================== 6 ====================
\chapter{Туршилт ба UI дүрслэл}
\section{Багшаас оноосон үндсэн дүрслэл}
Хичээлийн ерөнхий зааварт (Бие даалт 1, нийтлэг шаардлагын 4-р заалт) багшаас төсөлд оноосон үндсэн дүрслэл бүхий UI-ийг хөгжүүлж, бодит ажиллагаатай холбохыг шаардсан. GTK3 интерфэйс нь 9$\times$9 Sudoku хүснэгт, хэрэглэгчийн үйлдлийн товчлуурууд, шалгалтын үр дүн болон 11 шалгах урсгалын төлвийг харуулах бөгөөд PASS/FAIL, VALID/INVALID төлөв нь хэрэглэгчийн оруулсан бодит өгөгдлийг шалгасан үр дүнтэй холбогдсон.
\section{Туршилтын орчин}
Үйлдлийн систем: Windows 11; хөгжүүлэлтийн орчин: MSYS2 UCRT64; хөрвүүлэгч: GCC; урсгалын сан: POSIX Threads; график интерфэйс: GTK3; эх код: \repo. Нарийн хувилбарыг 5.2 хэсгийн хүснэгтээс харна уу.
\section{Юуг өөрчилж турших вэ}
Хэрэглэгч дараахыг өөрчилж туршиж болно: (1) 81 нүдний аль ч утга; (2) нүдийг хоосон үлдээх; (3) 1--9-өөс өөр тэмдэгт оруулах; (4) үр дүнг үзсэний дараа утгыг засаж дахин шалгах; (5) цэвэрлэх товчоор хүснэгтийг дахин эхлүүлэх. Ажиглах зүйл: хэсэг бүрийн PASS/FAIL, 11 урсгалын төлөв, эцсийн VALID/INVALID, мөн алдааны мэдэгдэл.
\section{Хоосон нүдтэй оролт}
Зөв хүснэгтийн (1,1) нүдийг хоосон үлдээсэн (Хүснэгт~\ref{fig:g1}-ийн тест 1). Програм шалгалт эхлүүлэхгүйгээр хоосон нүдний мөр, баганыг мэдэгдэх ёстой (Зураг~\ref{fig:empty-cell}).
\shot{gui-empty-cell.png}{Хоосон нүдтэй оролтын алдааны мэдээлэл}{fig:empty-cell}
\section{Буруу Sudoku}
Зөв хүснэгтийн (1,2) нүдийг 3-аас 5 болгож өөрчилсөн (тест 2). Мөр 1, багана 2, дэд хүснэгт 1 FAIL, эцсийн төлөв INVALID байх ёстой (Зураг~\ref{fig:invalid-sudoku}).
\shot{gui-invalid-sudoku.png}{Буруу Sudoku-ийн шалгалтын үр дүн}{fig:invalid-sudoku}
\section{Хэсэгчилсэн алдаа}
Хоёр тусдаа алдаатай хүснэгтэд (тест 3) алдаатай хэсгүүд FAIL, бусад нь PASS байж, мөр, багана, дэд хүснэгтийн төлөв тус тусдаа харагдах ёстой (Зураг~\ref{fig:partial-fail}).
\shot{gui-row-column-subgrid-fail.png}{Мөр, багана болон дэд хүснэгтийн шалгалтын үр дүн}{fig:partial-fail}
\section{Зөв Sudoku}
Зөв хүснэгтэд (тест 4) мөр, багана, есөн дэд хүснэгт бүгд PASS, эцсийн төлөв VALID байх ёстой (Зураг~\ref{fig:valid-sudoku}).
\shot{gui-valid-sudoku.png}{Зөв Sudoku-ийн шалгалтын үр дүн}{fig:valid-sudoku}
\section{Урсгалын төлөвийн самбар ба responsive интерфэйс}
Самбар нь 11 урсгал тус бүрийн төлвийг харуулна. Зөв хүснэгтэд бүгд PASS болсон байх ёстой (Зураг~\ref{fig:thread-status}).
\shot{gui-thread-status-pass.png}{11 урсгалын төлөв болон шалгалтын үр дүн}{fig:thread-status}

%==================== 7 ====================
\chapter{Туршилтын арга}
\section{Арга ба хүлээгдэх үр дүн}
Тест бүрд оролт, хүлээгдэж буй үр дүнг урьдчилан тодорхойлж, програмын бодит үр дүнтэй харьцуулсан. Дэд хүснэгтийг мөр-мөрөөр зүүнээс баруун тийш 1--9 гэж дугаарлана (дэд хүснэгт 1 нь зүүн дээд, 9 нь баруун доод). Нүдийг (мөр, багана) гэж заана. Тестийн оролтын хүснэгтүүдийг 7.2 хэсэгт бүрэн өгсөн тул өөр хүн ижил оролтоор давтаж шалгах боломжтой.
\begin{longtable}{P{0.06\textwidth}P{0.2\textwidth}P{0.32\textwidth}P{0.3\textwidth}}\caption{Туршилтын кейсүүд}\label{tab:tests}\\\toprule
\textbf{№}&\textbf{Тест}&\textbf{Оролт}&\textbf{Хүлээгдэж буй үр дүн}\\\midrule\endfirsthead
\toprule\textbf{№}&\textbf{Тест}&\textbf{Оролт}&\textbf{Хүлээгдэж буй үр дүн}\\\midrule\endhead
1&Хоосон нүд&Зөв хүснэгт, (1,1) хоосон&``1-р мөр, 1-р багана хоосон байна'' гэсэн алдааны мэдэгдэл\\
2&Буруу Sudoku&Зөв хүснэгт, (1,2)=5&Мөр 1, багана 2, дэд хүснэгт 1 FAIL; бусад PASS; INVALID\\
3&Хэсэгчилсэн алдаа&(1,1)=3 ба (9,9)=8&Мөр 1,9; багана 1,9; дэд хүснэгт 1,9 FAIL; бусад PASS; INVALID\\
4&Зөв Sudoku&Зөв хүснэгт&11 урсгал бүгд PASS; VALID\\
5&Зөвхөн мөр алдаатай&(1,1)ба(2,1)-ийг солисон&Мөр 1,2 FAIL; багана, дэд хүснэгт бүгд PASS; INVALID\\
6&Зөвхөн багана алдаатай&(1,1)ба(1,2)-ийг солисон&Багана 1,2 FAIL; мөр, дэд хүснэгт бүгд PASS; INVALID\\
7&Зөвхөн дэд хүснэгт алдаатай&Мөр 1 ба мөр 4-ийг солисон&Дэд хүснэгт 1--6 FAIL, 7--9 PASS; мөр, багана бүгд PASS; INVALID\\
8&Буруу тэмдэгт&Нүдэнд 0, 10, -1, a оруулах&0: ``1-р мөр, 1-р баганад 1--9 оруулна''; a, -: ижил алдаа; 10: нүд 1 тэмдэгтээр хязгаарлагдаж зөвхөн 1 үлдэнэ\\
9&Засаж дахин шалгах&Тест 2-ын дараа (1,2)-ийг 3 болгох&Дахин шалгахад бүх хэсэг PASS; VALID\\
10&Цэвэрлэх (reset)&Аль ч тестийн дараа Clear товч&Бүх нүд хоосон болж, статус ``Шалгаагүй'' болно; шалгалт явж байхад товч блоклогдоно\\
\bottomrule\end{longtable}
\section{Тестийн өгөгдөл}
Дараах хүснэгтүүдэд тест бүрийн оролтыг үзүүлэв. Эхний хүснэгт (Зөв) нь жишиг бөгөөд нүд хоосон байвал нүд хоосон гэж үзнэ; өөрчилсөн нүдийг тодруулав.
\begin{figure}[H]\centering
\begin{minipage}[t]{0.32\textwidth}\centering Тест 4 (Зөв)\par\sgrid{5&3&4&6&7&8&9&1&2}{6&7&2&1&9&5&3&4&8}{1&9&8&3&4&2&5&6&7}{8&5&9&7&6&1&4&2&3}{4&2&6&8&5&3&7&9&1}{7&1&3&9&2&4&8&5&6}{9&6&1&5&3&7&2&8&4}{2&8&7&4&1&9&6&3&5}{3&4&5&2&8&6&1&7&9}\end{minipage}\hfill
\begin{minipage}[t]{0.32\textwidth}\centering Тест 1 (Хоосон нүд)\par\sgrid{ &3&4&6&7&8&9&1&2}{6&7&2&1&9&5&3&4&8}{1&9&8&3&4&2&5&6&7}{8&5&9&7&6&1&4&2&3}{4&2&6&8&5&3&7&9&1}{7&1&3&9&2&4&8&5&6}{9&6&1&5&3&7&2&8&4}{2&8&7&4&1&9&6&3&5}{3&4&5&2&8&6&1&7&9}\end{minipage}\hfill
\begin{minipage}[t]{0.32\textwidth}\centering Тест 2 (Буруу)\par\sgrid{5&\textbf{5}&4&6&7&8&9&1&2}{6&7&2&1&9&5&3&4&8}{1&9&8&3&4&2&5&6&7}{8&5&9&7&6&1&4&2&3}{4&2&6&8&5&3&7&9&1}{7&1&3&9&2&4&8&5&6}{9&6&1&5&3&7&2&8&4}{2&8&7&4&1&9&6&3&5}{3&4&5&2&8&6&1&7&9}\end{minipage}
\caption{Тест 4, 1, 2-ийн оролт}\label{fig:g1}
\end{figure}
\begin{figure}[H]\centering
\begin{minipage}[t]{0.32\textwidth}\centering Тест 3 (Хэсэгчилсэн)\par\sgrid{\textbf{3}&3&4&6&7&8&9&1&2}{6&7&2&1&9&5&3&4&8}{1&9&8&3&4&2&5&6&7}{8&5&9&7&6&1&4&2&3}{4&2&6&8&5&3&7&9&1}{7&1&3&9&2&4&8&5&6}{9&6&1&5&3&7&2&8&4}{2&8&7&4&1&9&6&3&5}{3&4&5&2&8&6&1&7&\textbf{8}}\end{minipage}\hfill
\begin{minipage}[t]{0.32\textwidth}\centering Тест 5 (зөвхөн мөр)\par\sgrid{\textbf{6}&3&4&6&7&8&9&1&2}{\textbf{5}&7&2&1&9&5&3&4&8}{1&9&8&3&4&2&5&6&7}{8&5&9&7&6&1&4&2&3}{4&2&6&8&5&3&7&9&1}{7&1&3&9&2&4&8&5&6}{9&6&1&5&3&7&2&8&4}{2&8&7&4&1&9&6&3&5}{3&4&5&2&8&6&1&7&9}\end{minipage}\hfill
\begin{minipage}[t]{0.32\textwidth}\centering Тест 6 (зөвхөн багана)\par\sgrid{\textbf{3}&\textbf{5}&4&6&7&8&9&1&2}{6&7&2&1&9&5&3&4&8}{1&9&8&3&4&2&5&6&7}{8&5&9&7&6&1&4&2&3}{4&2&6&8&5&3&7&9&1}{7&1&3&9&2&4&8&5&6}{9&6&1&5&3&7&2&8&4}{2&8&7&4&1&9&6&3&5}{3&4&5&2&8&6&1&7&9}\end{minipage}
\caption{Тест 3, 5, 6-ын оролт}\label{fig:g2}
\end{figure}
\begin{figure}[H]\centering
\begin{minipage}[t]{0.4\textwidth}\centering Тест 7 (зөвхөн дэд хүснэгт)\par\sgrid{\textbf{8}&\textbf{5}&\textbf{9}&\textbf{7}&\textbf{6}&\textbf{1}&\textbf{4}&\textbf{2}&\textbf{3}}{6&7&2&1&9&5&3&4&8}{1&9&8&3&4&2&5&6&7}{\textbf{5}&\textbf{3}&\textbf{4}&\textbf{6}&\textbf{7}&\textbf{8}&\textbf{9}&\textbf{1}&\textbf{2}}{4&2&6&8&5&3&7&9&1}{7&1&3&9&2&4&8&5&6}{9&6&1&5&3&7&2&8&4}{2&8&7&4&1&9&6&3&5}{3&4&5&2&8&6&1&7&9}\end{minipage}
\caption{Тест 7-ийн оролт (мөр 1 ба 4 солигдсон)}\label{fig:g3}
\end{figure}
Тест 8--10-д тусгай хүснэгт хэрэггүй: тест 8-д дурын нүдэнд 0, 10, $-1$, \texttt{a} утга оруулж, тест 9-д тест 2-ын хүснэгтийн (1,2) нүдийг 3 болгож, тест 10-д аль ч хүснэгтийн дараа цэвэрлэх товч дарсан.

%==================== 8 ====================
\chapter{Үр дүн}
Тест бүрийн бодит үр дүнг хүлээгдэж буй үр дүнтэй харьцуулав. Нотолгоо болгон туршилтын үеийн дэлгэцийн зургуудыг ашигласан; тест 1--4-ийн зургийг 6-р бүлэгт үзүүлэв.
\begin{longtable}{P{0.06\textwidth}P{0.2\textwidth}P{0.4\textwidth}P{0.2\textwidth}}\caption{Үр дүнгийн нэгтгэл}\label{tab:results}\\\toprule
\textbf{№}&\textbf{Тест}&\textbf{Бодит үр дүн}&\textbf{Нотолгоо}\\\midrule\endfirsthead
\toprule\textbf{№}&\textbf{Тест}&\textbf{Бодит үр дүн}&\textbf{Нотолгоо}\\\midrule\endhead
1&Хоосон нүд&``Алдаа: 1-р мөр, 1-р багана хоосон байна.'' мэдэгдэл гарсан.&Зураг~\ref{fig:empty-cell}\\
2&Буруу Sudoku&INVALID төлөв гарч, зөрчилтэй хэсгүүд FAIL болсон.&Зураг~\ref{fig:invalid-sudoku}\\
3&Хэсэгчилсэн алдаа&Мөр, багана болон дэд хүснэгтийн PASS/FAIL төлвүүд тус тусдаа харагдсан.&Зураг~\ref{fig:partial-fail}\\
4&Зөв Sudoku&VALID мэдэгдэл гарч, мөр, багана, есөн дэд хүснэгт бүгд PASS болсон.&Зураг~\ref{fig:valid-sudoku}, \ref{fig:thread-status}\\
5&Зөвхөн мөр&Мөр 1, 2 FAIL; бүх багана, дэд хүснэгт PASS; нийт INVALID.&Дэлгэцийн зураг\\
6&Зөвхөн багана&Багана 1, 2 FAIL; бүх мөр, дэд хүснэгт PASS; нийт INVALID.&Дэлгэцийн зураг\\
7&Зөвхөн дэд хүснэгт&Дэд хүснэгт 1--6 FAIL, 7--9 PASS; бүх мөр, багана PASS; нийт INVALID.&Дэлгэцийн зураг\\
8&Буруу тэмдэгт&0 оруулахад ``Алдаа: 1-р мөр, 1-р баганад 1--9 оруулна.'' мэдэгдэл гарсан. 10, -1, a: нүд нэг тэмдэгтээр хязгаарлагддаг тул 10-аас зөвхөн эхний тэмдэгт үлдэж, ``-'' ба a нь цифр биш гэж алдаа өгнө (эх кодын дагуу).&Дэлгэцийн зураг (0)\\
9&Засаж дахин шалгах&Засварын дараа бүх хэсэг PASS, нийт VALID болсон.&Дэлгэцийн зураг\\
10&Цэвэрлэх&81 нүд хоосорч, статус ``Шалгаагүй'' болсон.&Дэлгэцийн зураг\\
\bottomrule\end{longtable}

%==================== 9 ====================
\chapter{Хэлэлцүүлэг}
Олон урсгал ашигласнаар Sudoku-ийн бие даасан хэсгүүдийг тус тусад нь шалгах боломжтой болсон. Тест 5--7 нь алдааг мөр, багана, дэд хүснэгтийн аль нэгэнд тусгаарлан үүсгэдэг тул PASS/FAIL үр дүн алдаа илэрсэн хэсгийг яг заадаг эсэхийг шалгана. GTK3 интерфэйс нь хүснэгт оруулах, шалгах, засварлан дахин шалгах боломж олгож, урсгалын төлөвийн самбараар 11 шалгалтын үр дүнг тус тусад нь харуулдаг.

\textbf{Хязгаарлалт.}
\begin{itemize}
\item Програм Sudoku-г бодож шийдэхгүй, зөвхөн өгөгдсөн шийдлийн зөв эсэхийг шалгана.
\item График хувилбарт GTK3 болон MSYS2 UCRT64 орчин шаардлагатай; Linux, macOS дээр туршаагүй.
\item Гүйцэтгэлийн хурдыг хэмжээгүй тул 11 урсгалын хурдны давуу талын тухай тоон дүгнэлт хийгээгүй. 9$\times$9 хүснэгт жижиг тул урсгал үүсгэх зардал нь шалгалтын зардлаас их байж болно.
\end{itemize}

%==================== 10 ====================
\chapter{Дүгнэлт}
Энэ төслөөр C хэл, POSIX Threads болон GTK3 ашигласан олон урсгалт Sudoku шалгагчийн консол болон график хувилбарыг боловсруулав. Шалгалтыг мөр, багана болон есөн 3$\times$3 дэд хүснэгтэд хуваарилж, нийт 11 урсгал ашигласан нь сурах бичгийн шаардлагатай нийцнэ. Хоосон нүд, буруу, хэсэгчилсэн алдаатай болон зөв Sudoku-г туршиж, үр дүнг screenshot-оор нотолсон. Цаашид оролтын баталгаажуулалтыг сайжруулах, Linux орчны заавар нэмэх, гүйцэтгэлийг хэмжих боломжтой.

%==================== ЛАВЛАГАА ====================
\chapter*{Ашигласан материал}\addcontentsline{toc}{chapter}{Ашигласан материал}
\begin{enumerate}
\item A. Silberschatz, P. B. Galvin, G. Gagne. \textit{Operating System Concepts}, 10th ed. Wiley, 2018. Chapter 4, Programming Project 5.
\item F.CSM302 Үйлдлийн систем, Бие даалт 1-ын заавар. Хичээлийн материал.
\item The Open Group. POSIX Threads API. \url{https://pubs.opengroup.org/onlinepubs/9699919799/}
\item GTK Project. GTK 3 Reference Manual. \url{https://docs.gtk.org/gtk3/}
\item GNU Compiler Collection. \url{https://gcc.gnu.org/}
\item Sudoku Validator төслийн эх код. \repo
\item Код, тайлбар, баримтжуулалт боловсруулахад ChatGPT, тайлангийн LaTeX засварт Claude ашигласан.
\end{enumerate}

%==================== ХАВСРАЛТ ====================
\appendix
\chapter{Ажлын төлөвлөгөө ба хуваарилалт}
\begin{table}[H]\centering\caption{Төслийн ерөнхий мэдээлэл}\begin{tabular}{P{0.28\textwidth}P{0.62\textwidth}}\toprule
\textbf{Талбар}&\textbf{Мэдээлэл}\\\midrule
Хуваарилалтын код&II-4-5\\
Төсөл&Sudoku Solution Validator\\
Repository&\repo\\
Хэл/сан/орчин&C, pthread, GTK3, GCC, MSYS2 UCRT64\\
Үндсэн дүрслэл&11 урсгалын төлвийн самбар, PASS/FAIL, VALID/INVALID\\\bottomrule\end{tabular}\end{table}
\begin{longtable}{P{0.3\textwidth}P{0.4\textwidth}P{0.2\textwidth}}
\caption{Багийн ажлын хуваарилалт}\\
\toprule
\textbf{Гишүүн} & \textbf{Хариуцсан ажил} & \textbf{Commit/file}\\
\midrule\endfirsthead
\toprule\textbf{Гишүүн} & \textbf{Хариуцсан ажил} & \textbf{Commit/file}\\\midrule\endhead
Д.~Алтан-Очир (B221940001) & Sudoku шалгах алгоритмын шаардлага, мөр болон багана шалгах логикийг хянах; алгоритмын тайлбар. & \texttt{sudoku\_core.c}\\
Б.~Аззаяа (B242270803) & GTK3 интерфэйсийн зохиомж, хэрэглэгчийн оролт болон товчлууруудын ажиллагааг шалгах. & \texttt{gui.c}\\
О.~Энхжаргал (B232230049) & Олон урсгалын зохицуулалт, \texttt{pthread} функцүүдийн тайлбар болон эх кодын хяналт. & \texttt{sudoku\_core.c/.h}\\
Н.~Мөнхбаяр (B242270058) & Туршилтын кейс бэлтгэх, зөв/буруу Sudoku-ийн үр дүнг шалгах, тестийн хүснэгт бүрдүүлэх. & тестийн оролт (тайлан 7.2)\\
Ц.~Тэргэл (B241940128) & Хоосон нүд болон алдаатай оролтын туршилт, дэлгэцийн зураг болон хэрэглэгчийн зааврын хяналт. & \texttt{docs/screenshots/}\\
А.~Найдан (B241930004) & GitHub repository, README/Manual-ийн уялдаа, тайланг нэгтгэх, LaTeX хувилбар, эцсийн шалгалт. & \texttt{README.md}, \texttt{main.c}, тайлан\\
\bottomrule
\end{longtable}
Commit-үүдийг хоёр нэрээр (\texttt{Naidan2006}, \texttt{Naidan}; ижил email) бүртгэсэн тул гишүүн тус бүрийн commit-ийн тоог тусад нь ялгах боломжгүй.
\end{document}